\documentclass[a4paper,11pt]{article}


% Please do not change the packages and commands listed here. You can of course add your favourite packages and commands.

%%% Pakete %%%

\usepackage{graphicx}
\usepackage{amsmath}
\usepackage[utf8]{inputenc}
\usepackage{enumerate}
\usepackage{amsthm}
\usepackage{amssymb}
\usepackage{paralist}
\usepackage{dsfont}
\usepackage{booktabs}
\usepackage{tabularx}
\usepackage{hyperref}
\usepackage{mathtools}
\usepackage{verbatim}
\usepackage{ dsfont }
\usepackage{ xcolor }
\usepackage{environ}
\usepackage{lipsum}


%%% Additional Commands %%%

%% These commands might mak it easier for you to write formal statements.

\newcommand{\N}{\ensuremath{\mathds{N}}} % natural numbers
\newcommand{\Z}{\ensuremath{\mathds{Z}}} % integers
\newcommand{\Q}{\ensuremath{\mathds{Q}}} % rational numbers
\newcommand{\R}{\ensuremath{\mathds{R}}} % real numbers

\newcommand{\states}{\ensuremath{\mathcal{S}}} % state space
\newcommand{\actions}{\ensuremath{\mathcal{A}}} % action space
\newcommand{\policy}{\ensuremath{\pi}} % a policy
\newcommand{\stateval}{\ensuremath{v}} % state value function
\newcommand{\actionval}{\ensuremath{q}} % action value function
\newcommand{\optsval}{\ensuremath{v_*}} % state value function
\newcommand{\optaval}{\ensuremath{q_*}} % action value function
\newcommand{\reward}{\ensuremath{\mathcal{R}}} % reward function
\newcommand{\transition}{\ensuremath{\mathcal{P}}} % transition function
\newcommand{\discount}{\ensuremath{\gamma}} % discount factor

% Wrap AI-assisted text in this to highlight it blue (required by course policy)
\newcommand{\aitxt}[1]{\textcolor{blue}{#1}}

\newtheorem{lemma}{Lemma}
\newtheorem{example}{Example}
\theoremstyle{definition}
\newtheorem{prob}{Exercise}
\newtheorem{pcode}{Pseudocode}




% -----------------------------
\begin{document}

%%% TITLE PAGE %%%
%   - Please only modify your name(s) and do not change anything else! %

\begin{titlepage}
\pagestyle{empty}
  \begin{minipage}[t]{0.6\textwidth}
    \begin{flushleft}
      \bf
      Reinforcement Learning and Decision Making Under Uncertainty\\
    \end{flushleft}
  \end{minipage}\hfill
  \begin{minipage}[t]{0.3\textwidth}
    \begin{flushright}
      \bf Spring 2022\\
      FistName, LastName
    \end{flushright}
  \end{minipage}

	\medskip

  \begin{center}
    {\Large\bf Project Report}

  {\bf Submission Deadline: 03 June 2022, 14:00}
    \Large
    \vspace{2cm}

    %%% ENTER THE NAMES OF GROUP MEMBERS HERE %%%
    %
    {Group:}\\
    {NAME 1}\\
    {NAME 2}

    \vspace{2cm}
    \begin{tabular}[b]{|l|r|}
    \hline
    Introduction, Related Work, Preliminaries &\qquad\qquad/\;\;5\\\hline
    Structure, Notation, Formal correctness &\qquad\qquad/\;\;5\\\hline
    Main Content and Explanations &\qquad\qquad/30\\\hline
    Discussion \& Conclusion&\qquad\qquad/10\\\hline
    $\Sigma$&\qquad\qquad/50\\\hline
    \end{tabular}
  \end{center}
\end{titlepage}
%%% END TITLE PAGE %%%




\newpage
\pagestyle{plain}
\pagenumbering{arabic}


\section{Introduction}
% What is the problem, why does it matter, and what does this report do?

\subsection*{Related Work}
% 2--3 relevant works and how they relate to this project. Cite with \cite{}.


\section{Preliminaries}

\subsection{MDP Formulation}
% Define the MDP tuple $(\states, \actions, \transition, \reward, \discount)$ for this problem.
% State the Bellman optimality equation and the definition of the optimal policy.


\section{Environment}

\subsection{Environment Description}
% Grid size, obstacle layout, start state, goal state(s). Include a figure.

\subsection{Implementation Details}
% State encoding, episode termination condition, stochasticity. Link to code repository.


\section{Algorithms}

\subsection{Q-Learning}
% Write the Q-learning update equation and describe the exploration strategy.

\subsection{SARSA}
%

\subsection{Deep Q-Network (DQN)}
% Write the DQN loss function. Mention replay buffer and target network.


\subsection{Hyperparameters}
% Table: learning rate, discount factor, exploration schedule, episodes, and DQN-specific params.


\section{Experiments}

\subsection{Hypotheses}
% List the specific questions the experiments aim to answer.

\subsection{Setup}
% Environment variant(s), number of runs, random seeds, evaluation metric, baselines.

\subsection{Results}
% Learning curves and/or tables. Figure: reward vs.\ episode.

\subsection{Discussion}
% Interpret the results. What explains the observed behaviour?


\section{Robustness}

\subsection{Robustness Analysis}
% Dimensions tested: stochasticity, unseen layouts, reward perturbations, hyperparameter sensitivity.

\subsection{Results}
% Table: performance under each condition. Discuss what the policy tolerates and where it breaks.


\section{Conclusion}

\subsection{Limitations}
% What does the approach not handle well?

\subsection{Future Work}
% Possible extensions to the project.


\section*{Use of External Aids}
% List every book, paper, URL, library, and AI tool used. Highlight AI-assisted text with \aitxt{}.


\bibliographystyle{alpha}
\bibliography{references}




\end{document}